\section{Related work}\label{sec:related}

\paragraph{CuTe and its specifications.}
CuTe grew out of Graphene~\cite{hagedorn2023} and is documented in the CUTLASS
repository~\cite{cutedocs}; Cecka's paper~\cite{cecka2026} is the authoritative specification. It
defines layouts with nested shapes and strides in an ``integer-semimodule'', states each operation
by post-conditions on an extended domain, and derives the stride and shape divisibility conditions
and, for distributivity over the modes of the inner layout, the condition that those modes be
mutually segregated, as sufficient conditions; it contains no theorems or proofs, and its
reference implementation PyCuTe~\cite{pycute} does not check the segregation condition. Our results need the stride domain to be a $\Z$-module, i.e.\ an
abelian group; the semimodule axioms of~\cite{cecka2026} omit distributivity, without which
subdividing a mode need not preserve the function.

\paragraph{Mathematical accounts.}
Shah~\cite{shah2024} gave the first rigorous treatment: admissibility conditions for complement
and composition with soundness proofs, the counterexample to unconditional distributivity that we
reuse in \cref{ex:shah}, and the observation that layout-expressible permutations come from a
groupoid. His single-mode composition rule is contained in our \cref{prop:monomial}(a,b); his
multi-mode
condition is pairwise and, as \cref{ex:three} shows, needs to be strengthened for three or more
modes.
Carlisle, Shah, Stern and VanKoughnett~\cite{carlisle2026} develop the theory of flat and nested
layouts with non-negative strides in detail---normal form of coalesce, compact layouts,
existence and uniqueness of complements---and introduce categories $\mathsf{Tuple}$ and
$\mathsf{Nest}$ whose morphisms encode \emph{tractable} layouts, with a graphical calculus and a
composition algorithm based on mutual refinement. Their categories and composition algorithm do
not cover layouts whose sorted strides fail the divisibility chain (e.g.\ $(4,8):(1,5)$) and give
no necessary condition for composability, and they work with non-negative integer strides. Our \cref{thm:spectrum} and
\cref{thm:complement} recover their normal-form and complement theorems by different arguments;
\cref{thm:exact} is exact and unrestricted, and \cref{thm:division} is exact for a compact inner
layout; \cref{sec:box} relates their category to ours.

\paragraph{Other layout formalisms.}
Triton's linear layouts~\cite{zhou2026} model layouts as $\F_2$-linear maps on bit vectors,
restricted to power-of-two shapes; \cref{sec:beyond} relates them to CuTe layouts with values in
$\F_2^k$. Bhaskaracharya et al.~\cite{bhaskaracharya2025} encode both CuTe layouts and linear
layouts as integer set relations and delegate the operations to the Integer Set Library; their
quasi-affine expressions are instances of our floor-sum form, and the chain condition of
\cref{thm:spectrum} characterises which of them are layouts. LEGO~\cite{tavakkoli2025} and
Axe~\cite{hou2026} are layout languages for code generators with different aims.

\paragraph{Swizzles.}
The CuTe documentation and source describe swizzles and composed layouts
operationally~\cite{cutedocs}. Cecka~\cite{cecka2026} obtains swizzles as layouts whose strides
are bit vectors added by XOR and composes one with a thread-value layout as an example; Triton's
linear layouts are objects of the same kind. We are not aware of an earlier statement of when the two
representations agree (\cref{cor:xorlinear}), of the decay criterion (\cref{thm:decay}), or of
the conditions of \cref{sec:swzops}.

\paragraph{Mechanised proofs.}
Liu's thesis~\cite{liu2026} formalises a fragment of the layout algebra (concatenation, coalesce,
composition) in Rocq and proves soundness theorems for it, together with a correctness result for
layout-permuted loop nests. As printed in~\cite[\S6.4.3]{liu2026}, the composition function
distributes over the modes of the inner layout and its soundness statement carries no
interference condition; by our reading of the printed definition it returns $(3,2):(2,3)$ on
\cref{ex:shah}. We have not examined the Rocq development itself.
Our Lean development (\cref{sec:lean}) is organised around the carry identity and proves the
composition formula together with its defect term, which is what makes the interference condition
explicit.

\paragraph{Number theory.}
The pairs of sets with $I\oplus J=\N$ were classified by de~Bruijn~\cite{debruijn1956}. The
finite case $I\oplus J=\{0,\dots,n-1\}$ is classical as well---such pairs are known as Krasner
factorisations, see~\cite{bodini2006}---and we use the formulation of Long~\cite{long1967}.
Bodini and Rivals~\cite{bodini2006} study the tiles of an interval algorithmically and note the uniqueness and involutivity of the ``dual'' that we use in
\cref{prop:unique-complement,cor:double}. That carries are $2$-cocycles is folklore; see Isaksen~\cite{isaksen2002}.
We are not aware of earlier uses of either in the context of tensor layouts.
